% Jawaban publik bernomor ganjil yang terhubung langsung dengan
% conditional_probability.tex: hanya latihan 13, 15, 17, 19, dan 21.
% Tidak ada jawaban instruktur terbatas yang ditambahkan. Latihan genap
% 14, 16, 18, 20, dan 22 dicatat sebagai kesenjangan penguasaan O001.
% Koreksi turunan: rujukan bagian pada jawaban 13(c) diperbaiki dari (a)
% menjadi (b-i), dan tata bahasa sumber pada jawaban 17(d) diperjelas.

% 13

\eocesol{(a)~Tidak, tetapi kita dapat menghitungnya jika A dan B saling independen.
(b-i)~0.21.
(b-ii)~0.79.
(b-iii)~0.3.
(c)~Tidak, karena 0.1 $\ne$ 0.21; nilai 0.21 dihitung berdasarkan independensi
pada bagian~(b-i).
(d)~0.143.}

% 15

\eocesol{(a)~Tidak, 0.18 responden termasuk dalam kombinasi ini.
(b)~$0.60 + 0.20 - 0.18 = 0.62$.
(c)~$0.18 / 0.20 = 0.9$.
(d)~$0.11 / 0.33 \approx 0.33$.
(e)~Tidak, karena jika saling independen, jawaban bagian~(c) dan~(d) akan sama.
(f)~$0.06 / 0.34 \approx 0.18$.}

% 17

\eocesol{(a)~Tidak. Ada 6~perempuan yang paling menyukai Five Guys Burgers.
(b)~$162 / 248 = 0.65$.
(c)~$181 / 252 = 0.72$.
(d)~Dengan asumsi independensi antara pemilihan pasangan kencan dan
preferensi hamburger--yang sepintas tampak wajar--diperoleh
$0.65 \times 0.72 = 0.468$.
(e)~$(252 + 6 - 1)/500 = 0.514$.}

% 19

\eocesol{(a) \\

\FigureFullPath[Diagram pohon dengan cabang utama “Dapat membuat diagram kotak?”
dan cabang sekunder “Lulus?”. Cabang utama “Dapat membuat diagram kotak” memiliki
dua kemungkinan, yaitu “Ya” dengan peluang 0.8 dan “Tidak” dengan peluang
0.2. Masing-masing cabang ini memiliki dua cabang sekunder. Cabang utama “Ya”
terbagi menjadi cabang “Ya” (untuk Lulus) dengan peluang bersyarat 0.86 dan
peluang akhir Ya-dan-Ya sebesar 0.688, serta cabang sekunder “Tidak” dengan
peluang bersyarat 0.14 dan peluang akhir Ya-dan-Tidak sebesar 0.112.
Cabang utama “Tidak” dari “Dapat membuat diagram kotak” memiliki cabang “Ya” dengan
peluang bersyarat 0.65 dan peluang akhir Tidak-dan-Ya sebesar 0.13, serta
cabang sekunder “Tidak” dengan peluang bersyarat 0.35 dan peluang akhir
Tidak-dan-Tidak sebesar 0.07.]{0.375}{ch_probability/figures/eoce/tree_drawing_box_plots/tree_drawing_box_plots}
(b)~0.84}

% 21

\eocesol{0.0714. Bahkan ketika pasien mendapat hasil tes positif untuk lupus,
peluang bahwa pasien tersebut benar-benar menderita lupus hanya 7,14\%.
House mungkin benar. \\

\FigureFullPath[Diagram pohon dengan cabang utama “Lupus” dan cabang sekunder
“Hasil” untuk tes lupus. Cabang utama “Lupus” memiliki dua kemungkinan, yaitu
“Ya” dengan peluang 0.02 dan “Tidak” dengan peluang 0.98. Masing-masing
cabang ini memiliki dua cabang sekunder. Cabang utama “Ya” terbagi menjadi cabang
“Ya” (untuk Hasil) dengan peluang bersyarat 0.98 dan peluang akhir
Ya-dan-Ya sebesar 0.0196, serta cabang sekunder “Tidak” dengan peluang
bersyarat 0.02 dan peluang akhir Ya-dan-Tidak sebesar 0.0004. Cabang utama
“Tidak” dari “Lupus” memiliki cabang sekunder “Ya” dengan peluang bersyarat
0.26 dan peluang akhir Tidak-dan-Ya sebesar 0.2548, serta cabang sekunder
“Tidak” dengan peluang bersyarat 0.74 dan peluang akhir Tidak-dan-Tidak
sebesar 0.7252.]{0.375}{ch_probability/figures/eoce/tree_lupus/tree_lupus.pdf}}
